\documentclass[11pt,a4paper]{amsart}
\usepackage[margin=2.6cm]{geometry}
\usepackage{amsmath,amssymb,amsthm}
\usepackage{stmaryrd}
\usepackage{graphicx}
\usepackage{booktabs}
\usepackage[colorlinks=true,linkcolor=blue,citecolor=blue,urlcolor=blue]{hyperref}
\usepackage{enumitem}

\newcommand{\Set}{\mathbf{Set}}
\newcommand{\Cat}{\mathbf{Cat}}
\newcommand{\Cof}{\mathbf{Cof}}
\newcommand{\DCont}{\mathbf{DCont}}
\newcommand{\Poly}{\mathbf{Poly}}
\newcommand{\cont}[2]{#1 \triangleleft #2}
\newcommand{\id}{\mathrm{id}}
\newcommand{\cod}{\mathrm{cod}}
\newcommand{\dom}{\mathrm{dom}}
\newcommand{\ob}{\mathrm{ob}\,}
\newcommand{\sh}{\sharp}
\newcommand{\co}{\mathrel{\text{\reflectbox{$\rightharpoonup$}}}}
\DeclareMathOperator{\Hom}{Hom}

\theoremstyle{plain}
\newtheorem{theorem}{Theorem}[section]
\newtheorem{proposition}[theorem]{Proposition}
\newtheorem{lemma}[theorem]{Lemma}
\newtheorem{corollary}[theorem]{Corollary}
\theoremstyle{definition}
\newtheorem{definition}[theorem]{Definition}
\newtheorem{example}[theorem]{Example}
\theoremstyle{remark}
\newtheorem{remark}[theorem]{Remark}

\title[The morphisms of directed containers are cofunctors]{The morphisms of directed containers are cofunctors,\\ not functors}
\author{MacBeth}
\address{Kodamai}
\date{10 June 2026}

\begin{document}

\begin{abstract}
Directed containers are equivalent, at the level of objects, to small categories:
this dictionary underlies their use as a syntax for composition in supply chains,
agent orchestration, and bidirectional transformation. It is tempting to extend the
dictionary to morphisms by declaring that a morphism of directed containers is a
functor. We show this is false on grounds of variance. The position component of a
container morphism pulls positions \emph{backward}, $f^\sharp_s\colon P'(fs)\to P(s)$,
whereas a functor pushes morphisms \emph{forward}; the two variances are opposite.
Writing out the compatibility conditions of a \emph{directed}-container morphism, we
find they are exactly the cofunctor (retrofunctor) laws on the associated categories.
The resulting assignment is an isomorphism of categories $\DCont\cong\Cof$, where
$\Cof$ is the category of small categories and cofunctors --- not $\DCont\simeq\Cat$,
which fails already on a one-object source. The opposite (covariant) variance recovers
functors exactly, so functors and cofunctors are the two variances of a single notion
of structured map. This is not a defect but the point: the contravariant half is
precisely the \emph{Put} of a delta lens, so a directed-container morphism is the
update-propagating component of a bidirectional transformation. All counts are
verified combinatorially over $36$ ordered pairs of small categories.
\end{abstract}

\maketitle

\section{Introduction}

A directed container \cite{ACU14,AU16} is a container $\cont SP$ equipped with the
data of an indexed monoid: a root $\mathsf o(s)\in P(s)$, a sub-shape operation
$s\downarrow p$, and a shift $p\oplus q$ that compose positions associatively with the
root as unit. The motivating theorem is that this data is \emph{exactly} the data of a
small category: shapes are objects, positions out of $s$ are morphisms out of $s$, the
root is the identity, the shift is composition, and $s\downarrow p$ is the codomain
of $p$ \cite{AU16}. This object-level dictionary
\[
  \text{directed containers}\ \cong\ \text{small categories}
\]
is what lets one read a category as a coalgebraic, position-based structure --- the
form in which composition is implemented in applications: a supply chain as composable
sub-shapes, an agent hierarchy as composable positions, a state space as composable
updates.

A dictionary on objects invites a dictionary on morphisms. The roadmap slogan --- and a
natural first guess --- is that a morphism of directed containers should be a
\emph{functor} of the associated categories, upgrading the bijection to an equivalence
$\DCont\simeq\Cat$. This note shows the slogan is wrong, identifies the correct
statement, and argues that the correction is exactly what makes the bidirectional
applications work.

\paragraph{The variance problem.} A morphism of containers $(\cont SP)\to(\cont{S'}{P'})$
is a pair $(f,f^\sharp)$ with $f\colon S\to S'$ \emph{forward} on shapes but
$f^\sharp_s\colon P'(fs)\to P(s)$ \emph{backward} on positions: it takes a position out
of the target shape $fs$ and returns a position out of the source shape $s$. Under the
dictionary $f^\sharp_s$ sends a morphism out of $fs$ to a morphism out of $s$. A functor,
by contrast, sends a morphism out of $s$ to a morphism out of $Fs$, a map
$P(s)\to P'(Fs)$ running the \emph{other way}. The variances are opposite, so ``a
directed-container morphism is a functor'' is not merely unproved --- it is ill-typed.

\paragraph{The result.} The maps with this backward variance, subject to the directed
structure, have a name. A \emph{cofunctor} (retrofunctor) \cite{Aguiar,Clarke20}
$\varphi\colon\mathcal C\co\mathcal D$ is an object map together with a lifting operation
that takes an object $c$ and a morphism out of $\varphi_0 c$ and returns a morphism out
of $c$, subject to three laws. Writing the compatibility conditions
(M0)--(M2) of a directed-container morphism and reading them through the dictionary, we
find they \emph{are} the cofunctor laws (C0)--(C2), term for term. The assignment is a
bijection on every hom-set and respects composition, and it is a bijection on objects;
hence:
\[
  \boxed{\ \DCont\ \cong\ \Cof\ }
\]
an isomorphism of categories, where $\Cof$ is the category of small categories and
cofunctors (Theorem~\ref{thm:main}). The would-be statement $\DCont\simeq\Cat$ is false:
already over a one-object source there are strictly fewer cofunctors than functors
(Proposition~\ref{prop:neq}), so no full functor $\DCont\to\Cat$ can agree with the
object dictionary.

\paragraph{The method, and why it is not a defect.} The proof is a translation: the
data $(f,f^\sharp)$ and a cofunctor $(\varphi_0,\varphi^\uparrow)$ are literally the same
data, and (M0)--(M2) match (C0)--(C2) under the substitution
``position out of $s$ $=$ morphism out of $s$'' and ``$\oplus$ $=$ reverse
composition''. The covariant variant --- a position map $P(s)\to P'(fs)$ subject to the
mirror-image conditions --- recovers functors exactly (Proposition~\ref{prop:cov}). So a
directed structure carries \emph{two} natural notions of map, of opposite variance:
covariant gives functors, contravariant gives cofunctors, and only the latter is a
morphism of the underlying container.

This is the useful half. A \emph{delta lens} \cite{JR13,Clarke20,Clarke22} --- the
categorical model of a bidirectional transformation --- is precisely a \emph{Get}
functor together with a \emph{Put} cofunctor sharing an object map, related by a section
condition. Under $\DCont\cong\Cof$, the Put cofunctor \emph{is} a directed-container
morphism. So the contravariant variance flagged as a problem above is exactly the
component that propagates an update from a view back to its source
(Corollary~\ref{cor:lens}). Reading a directed-container map as a functor would silently
discard the update semantics --- and, as we show, even change the count of maps.

\paragraph{Outline.} Section~\ref{sec:prelim} recalls directed containers, the object
dictionary, and cofunctors. Section~\ref{sec:dict} writes the morphism conditions and
proves $\DCont\cong\Cof$. Section~\ref{sec:functors} places functors as the opposite
variance and shows the dictionary does not extend to $\Cat$. Section~\ref{sec:lens}
identifies a directed-container morphism with the Put of a delta lens.
Section~\ref{sec:verify} reports the combinatorial verification.
Section~\ref{sec:related} places the result among its 2026 neighbours, and
Section~\ref{sec:conclusion} discusses consequences and what remains.

\section{Directed containers and cofunctors}\label{sec:prelim}

We recall only what the proof uses, following \cite{ACU14,AU16} for containers and
\cite{Aguiar,Clarke20} for cofunctors.

\subsection{Directed containers and the object dictionary}

A \emph{container} $\cont SP$ is a set $S$ of \emph{shapes} together with a family
$P\colon S\to\Set$ of \emph{positions}. A \emph{directed container} structure on
$\cont SP$ is a root $\mathsf o(s)\in P(s)$, a sub-shape map $s\downarrow p\in S$ for
$p\in P(s)$, and a shift $p\oplus q\in P(s)$ for $p\in P(s)$, $q\in P(s\downarrow p)$,
satisfying the indexed-monoid laws (root is a two-sided unit for $\oplus$; $\oplus$ is
associative; $\downarrow$ acts compatibly). We refer to \cite{AU16} for the precise
axioms; we will use them only through the dictionary below.

\begin{theorem}[Object dictionary \cite{AU16}]\label{thm:obj}
Directed-container structures on $\cont SP$ are in bijection with small categories
$\mathcal C$ with object set $S$ and, for each $s$, $P(s)=\coprod_{s'}\mathcal C(s,s')$
the set of morphisms out of $s$, under
\[
  \ob\mathcal C=S,\quad
  \mathcal C(s,s')=\{p\in P(s): s\downarrow p=s'\},\quad
  \id_s=\mathsf o(s),\quad
  q\circ p = p\oplus q .
\]
In particular $s\downarrow p=\cod(p)$, and every small category arises from a unique
directed container.
\end{theorem}

Throughout, $\mathcal C(D)$ denotes the category of a directed container $D$. We will
move freely across the dictionary: a position $p\in P(s)$ \emph{is} a morphism out of
$s$, with codomain $s\downarrow p$; the root \emph{is} the identity; and $\oplus$ is
composition in reverse order, $p\oplus q = q\circ p$.

\subsection{Cofunctors}

\begin{definition}[Cofunctor \cite{Aguiar,Clarke20}]\label{def:cof}
A \emph{cofunctor} $\varphi\colon\mathcal C\co\mathcal D$ between small categories
consists of
\begin{itemize}[nosep]
  \item an object map $\varphi_0\colon\ob\mathcal C\to\ob\mathcal D$;
  \item a \emph{lift}: for each object $c$ of $\mathcal C$ and each morphism
        $h\colon\varphi_0 c\to d$ of $\mathcal D$, a morphism
        $\varphi^\uparrow(c,h)\colon c\to c'$ of $\mathcal C$,
\end{itemize}
subject to:
\begin{align}
  \text{(C0)}\quad & \varphi_0\big(\cod\varphi^\uparrow(c,h)\big)=\cod(h)
        && \text{(codomain)}\notag\\
  \text{(C1)}\quad & \varphi^\uparrow(c,\id_{\varphi_0 c})=\id_c
        && \text{(unit)}\notag\\
  \text{(C2)}\quad & \varphi^\uparrow(c,h'\circ h)
                   =\varphi^\uparrow(c',h')\circ\varphi^\uparrow(c,h),
        \quad c'=\cod\varphi^\uparrow(c,h)
        && \text{(composition)}\notag
\end{align}
In (C2), $h\colon\varphi_0 c\to d$ and $h'\colon d\to d''$; the term
$\varphi^\uparrow(c',h')$ typechecks because $\varphi_0 c'=d=\dom(h')$ by (C0).
Cofunctors compose --- object maps forward, lifts backward,
$(\psi\varphi)^\uparrow(c,m)=\varphi^\uparrow\big(c,\psi^\uparrow(\varphi_0 c,m)\big)$
--- with identities $(\id,\id^\uparrow)$, forming a category $\Cof$.
\end{definition}

A cofunctor is \emph{not} a contravariant functor \cite[p.\,188]{Clarke20}: it does not
act on all of $\mathcal C$'s morphisms, only lifting morphisms presented at the
$\varphi_0$-image of an object. Equivalently \cite{Clarke20}, a cofunctor is a span
$\mathcal C\xleftarrow{\ \varphi\ }X\xrightarrow{\ \bar\varphi\ }\mathcal D$ with
$\varphi$ bijective-on-objects and $\bar\varphi$ a discrete opfibration; cofunctors are
the comonoid morphisms in $(\Poly,\triangleleft)$ \cite{AU16,SS}. We will use only
the elementary form of Definition~\ref{def:cof}.

\section{The morphism dictionary: \texorpdfstring{$\DCont\cong\Cof$}{DCont = Cof}}\label{sec:dict}

We now write the conditions that make a container morphism a morphism of \emph{directed}
containers, and read them as the cofunctor laws.

\begin{definition}[Directed-container morphism]\label{def:dcm}
A \emph{morphism of directed containers} $D\to D'$ is a container morphism
$(f,f^\sharp)$, $f\colon S\to S'$ and $f^\sharp_s\colon P'(fs)\to P(s)$, satisfying for
all $s\in S$, $h\in P'(fs)$ and $k\in P'\big((fs)\downarrow' h\big)$:
\begin{align}
  \text{(M0)}\quad & f\big(s\downarrow f^\sharp_s(h)\big) = (fs)\downarrow' h
        && (\downarrow\text{-compatibility})\notag\\
  \text{(M1)}\quad & f^\sharp_s\big(\mathsf o'(fs)\big) = \mathsf o(s)
        && \text{(root)}\notag\\
  \text{(M2)}\quad & f^\sharp_s(h\oplus' k)
        = f^\sharp_s(h)\ \oplus\ f^\sharp_{\,s\downarrow f^\sharp_s(h)}(k)
        && \text{(shift homomorphism)}\notag
\end{align}
\end{definition}

\begin{remark}[Why (M2) typechecks]
Put $p:=f^\sharp_s(h)\in P(s)$. The left side applies $f^\sharp_s$ to
$h\oplus' k\in P'(fs)$. On the right, (M0) gives $(fs)\downarrow' h=f(s\downarrow p)$, so
$k\in P'\big(f(s\downarrow p)\big)$ and $f^\sharp_{s\downarrow p}(k)\in P(s\downarrow p)$;
hence $p\oplus f^\sharp_{s\downarrow p}(k)\in P(s)$ is defined. The shift condition is
well-formed precisely \emph{because} of the codomain condition --- exactly as (C2)
needs (C0).
\end{remark}

\begin{lemma}[Hom-set bijection]\label{lem:hom}
Let $D,D'$ be directed containers with categories $\mathcal C=\mathcal C(D)$,
$\mathcal C'=\mathcal C(D')$. The assignment
\[
  (f,f^\sharp)\ \longmapsto\
  \big(\varphi_0:=f,\ \ \varphi^\uparrow(s,h):=f^\sharp_s(h)\big)
\]
is a bijection between morphisms $D\to D'$ of directed containers
(Definition~\ref{def:dcm}) and cofunctors $\mathcal C\co\mathcal C'$
(Definition~\ref{def:cof}).
\end{lemma}

\begin{proof}
The key observation is that there is nothing to construct: under the dictionary the two
sides are the \emph{same data}. A position $h\in P'(fs)$ is a morphism out of $fs$ in
$\mathcal C'$ with codomain $(fs)\downarrow' h$, and $f^\sharp_s(h)\in P(s)$ is a
morphism out of $s$ in $\mathcal C$ with codomain $s\downarrow f^\sharp_s(h)$. So $f$ is
an object map and $f^\sharp$ is a lift; what must be checked is that the conditions
correspond. We match them term by term.

\emph{(M0)$\Leftrightarrow$(C0).} (M0) reads
$f\big(s\downarrow f^\sharp_s(h)\big)=(fs)\downarrow' h$. The left side is
$\varphi_0\big(\cod\varphi^\uparrow(s,h)\big)$; the right side is $\cod(h)$ in
$\mathcal C'$. This is (C0).

\emph{(M1)$\Leftrightarrow$(C1).} The identity at $fs$ is $\mathsf o'(fs)=\id_{fs}$ and
the identity at $s$ is $\mathsf o(s)=\id_s$, so $f^\sharp_s(\mathsf o'(fs))=\mathsf o(s)$
is $\varphi^\uparrow(s,\id_{\varphi_0 s})=\id_s$. This is (C1).

\emph{(M2)$\Leftrightarrow$(C2).} Translate $\oplus$ into composition by
$x\oplus y=y\circ x$. With $p=f^\sharp_s(h)$ and $s'=s\downarrow p$, the right side of
(M2) is
\[
  p\oplus f^\sharp_{s'}(k)=f^\sharp_{s'}(k)\circ f^\sharp_s(h)
  =\varphi^\uparrow(s',k)\circ\varphi^\uparrow(s,h),
\]
and the left side is
$f^\sharp_s(h\oplus' k)=f^\sharp_s(k\circ' h)=\varphi^\uparrow(s,k\circ' h)$. So (M2)
reads $\varphi^\uparrow(s,k\circ' h)=\varphi^\uparrow(s',k)\circ\varphi^\uparrow(s,h)$
with $s'=\cod\varphi^\uparrow(s,h)$, which is (C2).

The three conditions correspond under the identity map on data, so the assignment is a
bijection on hom-sets, with inverse the same identification read backward.
\end{proof}

\begin{lemma}[Functoriality]\label{lem:fun}
$\mathcal C(-)$ with the assignment of Lemma~\ref{lem:hom} is a functor
$\Phi\colon\DCont\to\Cof$.
\end{lemma}

\begin{proof}
The identity container morphism $(\id_S,\id)$ maps to $(\id,\id^\uparrow)$, the identity
cofunctor. For composition, let $(f,f^\sharp)\colon D\to D'$ and
$(g,g^\sharp)\colon D'\to D''$. The container composite has object map $g\circ f$ and
position map
\[
  (g\circ f)^\sharp_s = f^\sharp_s\circ g^\sharp_{fs}\colon
  P''(gfs)\xrightarrow{\ g^\sharp_{fs}\ }P'(fs)\xrightarrow{\ f^\sharp_s\ }P(s),
\]
positions pulling back through both stages. Under $\Phi$ this is the cofunctor with
object map $g_0\circ f_0$ and lift
$m\mapsto f^\sharp_s\big(g^\sharp_{fs}(m)\big)
=\varphi^\uparrow\big(s,\psi^\uparrow(\varphi_0 s,m)\big)$, which is the composite
cofunctor $\psi\varphi$ of Definition~\ref{def:cof}. So $\Phi$ preserves composition.
\end{proof}

\begin{theorem}[The morphism dictionary]\label{thm:main}
$\Phi\colon\DCont\to\Cof$ is an isomorphism of categories,
\[
  \DCont\ \cong\ \Cof .
\]
\end{theorem}

\begin{proof}
$\Phi$ is fully faithful by Lemma~\ref{lem:hom} (a bijection on every hom-set) and a
functor by Lemma~\ref{lem:fun}. On objects $D\mapsto\mathcal C(D)$ is a bijection by
Theorem~\ref{thm:obj}. A bijective-on-objects, fully faithful functor is an isomorphism
of categories: the object bijection and the hom bijections assemble into a strict
two-sided inverse.
\end{proof}

\section{Functors are the opposite variance}\label{sec:functors}

Theorem~\ref{thm:main} says the maps native to directed containers are cofunctors. Where,
then, do functors sit? At the opposite variance --- and \emph{not} as a refinement of
the same dictionary.

\begin{proposition}[Functors $=$ covariant-position morphisms]\label{prop:cov}
Call an \emph{op-variance morphism} $D\to D'$ a pair $(f,\phi)$ with $f\colon S\to S'$
and a \emph{covariant} position map $\phi_s\colon P(s)\to P'(fs)$ satisfying
\textup{(N0)} $f(s\downarrow p)=(fs)\downarrow'\phi_s(p)$,
\textup{(N1)} $\phi_s(\mathsf o(s))=\mathsf o'(fs)$, and
\textup{(N2)} $\phi_s(p\oplus q)=\phi_s(p)\oplus'\phi_{s\downarrow p}(q)$.
Then $(f,\phi)\mapsto F$, with $F=f$ on objects and $F(p)=\phi_s(p)$ on a morphism $p$
out of $s$, is a bijection onto functors $\mathcal C(D)\to\mathcal C(D')$.
\end{proposition}

\begin{proof}
A functor is exactly an object map $f$, a map of out-sets
$\phi_s=F|_{P(s)}\colon P(s)\to P'(fs)$ for each $s$, preserving codomains (N0),
identities (N1), and composition (N2: $F(q\circ p)=F(q)\circ F(p)$, i.e.\
$F(p\oplus q)=\phi_s(p)\oplus'\phi_{s\downarrow p}(q)$). Each datum and law matches under
the identity on data.
\end{proof}

So a directed structure carries two notions of structured map, of opposite variance:
\[
  \begin{aligned}
  \text{covariant positions }P(s)\to P'(fs)\ &\leadsto\ \textbf{functors},\\
  \text{contravariant positions }P'(fs)\to P(s)\ &\leadsto\ \textbf{cofunctors},
  \end{aligned}
\]
and only the contravariant one is a morphism of the underlying container. This is not a
matter of which is ``better''; it is that the container forgets to the second.

The two are genuinely different, even in cardinality.

\begin{proposition}[The dictionary does not extend to $\Cat$]\label{prop:neq}
Let $\Phi_0\colon\ob\DCont\to\ob\Cat$, $D\mapsto\mathcal C(D)$, be the object dictionary
of Theorem~\ref{thm:obj}. There is no functor $\DCont\to\Cat$ that agrees with $\Phi_0$
on objects and is full. In particular the object dictionary does not extend to an
equivalence $\DCont\simeq\Cat$.
\end{proposition}

\begin{proof}
Let $D_{\mathbf 1},D_{\mathbf 2}$ be the directed containers of the terminal category
$\mathbf 1$ and of $\mathbf 2=\{0\xrightarrow{u}1\}$. By Theorem~\ref{thm:main},
$|\DCont(D_{\mathbf 1},D_{\mathbf 2})|=|\Cof(\mathbf 1,\mathbf 2)|$, and we count the
latter directly. A cofunctor $\mathbf 1\co\mathbf 2$ picks $\varphi_0(\ast)\in\{0,1\}$
and must lift every morphism out of $\varphi_0(\ast)$ to the unique morphism $\id_\ast$
out of $\ast$, with (C0) forcing $\varphi_0(\cod(\text{lift}))=\cod$. Choosing
$\varphi_0(\ast)=0$ fails: $u\colon 0\to 1$ would lift to $\id_\ast$, but then (C0)
demands $\varphi_0(\ast)=\cod(u)=1\neq 0$. Choosing $\varphi_0(\ast)=1$ works: the only
morphism out of $1$ is $\id_1$, lifting to $\id_\ast$. So
$|\Cof(\mathbf 1,\mathbf 2)|=1$, whence $|\DCont(D_{\mathbf 1},D_{\mathbf 2})|=1$. But
$|\Cat(\mathbf 1,\mathbf 2)|=2$ (a functor $\mathbf 1\to\mathbf 2$ chooses either
object). A functor $G\colon\DCont\to\Cat$ equal to $\Phi_0$ on objects would restrict to
a map from a $1$-element hom-set onto a $2$-element hom-set, which cannot be surjective;
so $G$ is not full.
\end{proof}

\begin{remark}[$\Cof$ and $\Cat$ are incomparable, not nested]
The failure is not that $\Cof$ sits inside $\Cat$ with fewer maps. The comparison is not
even monotone over common objects: for the $3$-chain $\mathbf 3$,
$|\Cat(\mathbf 3,\mathbf 2)|=4<5=|\Cof(\mathbf 3,\mathbf 2)|$, while
$|\Cat(\mathbf 2,\mathbf 2)|=3>2=|\Cof(\mathbf 2,\mathbf 2)|$
(Table~\ref{tab:counts}). Neither hom-set contains the other. In $\Poly$ this is the
statement that functors and cofunctors occupy different hom-objects: cofunctors are the
comonoid morphisms in $(\Poly,\triangleleft)$, whereas functors are carried by
bicomodules \cite{AU16,SS}. We do not claim a separation of $\Cof$ and $\Cat$ as
abstract categories; Proposition~\ref{prop:neq} settles the question the dictionary
actually poses.
\end{remark}

\section{The cofunctor is the Put of a delta lens}\label{sec:lens}

The contravariant variance is the useful one. A \emph{delta lens} \cite{JR13,Clarke20}
is the categorical model of a bidirectional transformation between a source category $A$
and a view category $B$: it propagates a chosen update on the view back to a coherent
update on the source.

\begin{definition}[Delta lens \cite{Clarke22}]\label{def:lens}
A \emph{delta lens} $(f,\varphi)\colon A\rightleftarrows B$ consists of a functor
$f\colon A\to B$ (the \emph{Get}) and, for each object $a\in A$ and each morphism
$u\colon fa\to b$ in $B$, a morphism $\varphi(a,u)\colon a\to a'$ in $A$ (the
\emph{Put}/lift), satisfying
\begin{align}
  \text{(PutGet)}\quad & f\cdot\varphi(a,u)=u, \notag\\
  \text{(PutId)}\quad  & \varphi(a,\id_{fa})=\id_a, \notag\\
  \text{(PutPut)}\quad & \varphi(a,v\circ u)=\varphi(a',v)\circ\varphi(a,u),
        \quad a'=\cod\varphi(a,u). \notag
\end{align}
\end{definition}

The Put data $\varphi$ alone --- an object map (here $f$ on objects) plus a lift
satisfying the codomain, unit and composition laws --- is exactly a cofunctor
$A\co B$ \cite[Def.\,3]{Clarke22} (object map $f_0\colon A_0\to B_0$, lifts landing in
the source $A$, matching the convention of Definition~\ref{def:cof}). The Get $f$ supplies the functor; (PutGet) is the
\emph{section condition} relating the two, asserting that applying Get to a lifted
update returns the original update $u$. Clarke makes the decomposition precise: an
(internal) lens is a functor together with a cofunctor on a \emph{shared object map},
related by (PutGet); equivalently a commuting triangle
$A\xleftarrow{\varphi}\Lambda\xrightarrow{\bar\varphi}B$ with $\varphi$
identity-on-objects and $\bar\varphi$ a discrete opfibration, and $f\circ\varphi=\bar\varphi$
\cite[Cor.\,20]{Clarke20}. Indeed $f\colon\mathrm{DLens}(B)\to\Cof(B)$ forgetting the Get
is comonadic \cite[Thm.\,9]{Clarke22}.

\begin{remark}[A correction worth flagging]
The section condition is (PutGet), $f\cdot\varphi(a,u)=u$, \emph{not} ``$f\circ\varphi=\id$''.
Applying Get to a lifted update returns the original update, not an identity; the latter
holds only when read as a section of the fibre of updates over each object. We state this
because the loose form recurs.
\end{remark}

Combining Theorem~\ref{thm:main} with the decomposition gives the application hook.

\begin{corollary}[A directed-container morphism is the Put of a lens]\label{cor:lens}
Under the isomorphism $\DCont\cong\Cof$ of Theorem~\ref{thm:main}, the Put cofunctor of a
delta lens $(f,\varphi)\colon A\rightleftarrows B$ corresponds to a morphism of directed
containers $D_A\to D_B$ between the directed containers of $A$ and $B$. The covariant
component recovered by Proposition~\ref{prop:cov} is the Get functor. Hence a delta
lens is a directed-container morphism (the Put) together with a functor on the same
shapes (the Get), tied by (PutGet).
\end{corollary}

\begin{proof}
Immediate from Theorem~\ref{thm:main} and Definition~\ref{def:lens}: the Put $\varphi$ is
a cofunctor $A\co B$ \cite{Clarke22}, which $\Phi^{-1}$ carries to a directed-container
morphism $D_A\to D_B$; Proposition~\ref{prop:cov} identifies the Get with the covariant
morphism on the same object map.
\end{proof}

This is why the variance matters in applications. Whenever a construction is transported
along the equivalence chain --- composing supply-chain stages, propagating a change
through an agent hierarchy, updating a derived view of a state space --- the maps that
carry update semantics are the cofunctors, i.e.\ the directed-container morphisms.
Reading them as functors keeps the Get and discards the Put.

\section{Computational verification}\label{sec:verify}

We verified the correspondence combinatorially. A script represents a small category
concretely and enumerates, between two categories $C,D$: functors $C\to D$; cofunctors
$C\co D$ via (C0)--(C2); directed-container morphisms via (M0)--(M2) computed
\emph{directly} from the container data $(\mathsf o,\downarrow,\oplus)$; and op-variance
morphisms via (N0)--(N2). Over all $36$ ordered pairs from
$\{\mathbf 1,\mathbf 2,\mathbf 3,\mathbb Z/2,\mathbb Z/3,M\}$, with $M$ the two-element
idempotent monoid $\{1,t\}$, $t^2=t$:
\[
  \begin{aligned}
  \#\{\text{DCont-morphisms}\}&=\#\{\text{cofunctors}\}\quad\text{in all }36\text{ cases},\\
  \#\{\text{op-variance morphisms}\}&=\#\{\text{functors}\}\quad\text{in all }36\text{ cases},
  \end{aligned}
\]
confirming Lemma~\ref{lem:hom} and Proposition~\ref{prop:cov}; while
$\#\{\text{functors}\}\neq\#\{\text{cofunctors}\}$ in $17$ of the $36$ cases, confirming
Proposition~\ref{prop:neq}. A representative excerpt:

\begin{table}[h]\centering
\begin{tabular}{lrrr}
\toprule
$C\to D$ & \#functors & \#cofunctors & \#DCont-morphisms\\
\midrule
$\mathbf 1\to\mathbf 2$ & $2$ & $1$ & $1$\\
$\mathbf 2\to\mathbf 2$ & $3$ & $2$ & $2$\\
$\mathbf 2\to\mathbb Z/3$ & $3$ & $1$ & $1$\\
$\mathbf 3\to\mathbf 2$ & $4$ & $5$ & $5$\\
$\mathbf 3\to\mathbf 3$ & $10$ & $6$ & $6$\\
$\mathbb Z/3\to\mathbb Z/3$ & $3$ & $3$ & $3$\\
$\mathbf 3\to M$ & $4$ & $5$ & $5$\\
\bottomrule
\end{tabular}
\caption{Functors vs.\ cofunctors vs.\ directed-container morphisms.
\#cofunctors $=$ \#DCont-morphisms throughout; \#functors differs. The rows
$\mathbf 3\to\mathbf 2$ and $\mathbf 2\to\mathbf 2$ show that neither of $\Cat,\Cof$
dominates the other.}
\label{tab:counts}
\end{table}

\section{Related work}\label{sec:related}

The object dictionary of Theorem~\ref{thm:obj} and the retrofunctor reading of its morphisms
are, by 2026, common infrastructure. Two recent papers build over that infrastructure in
directions orthogonal to ours, and locate the present note among its neighbours without
displacing it. Both take the equivalence between retrofunctors and polynomial comonoids
\cite{AU16,Clarke20} as \emph{black-box} prior art --- exactly the identification that
Theorem~\ref{thm:main} unpacks at the level of container data.

Garner, Renata, and Wu \cite{GRW26} construct a contravariant idempotent adjunction between
ranked monads on $\Set$ and internal categories-with-retrofunctors in the category of locales;
its fixed points are the hyperaffine-unary monads on one side and the ample localic categories
on the other (their ``Stone duality for monads''). Their starting observation is ours: ``the
notion of morphism we need is not functors, but rather retrofunctors'' \cite[\S1]{GRW26}, and
they read a retrofunctor $\mathcal C\co\mathcal D$ as a \emph{simulation of transition systems}
--- the transition-system avatar of the update-propagating Put isolated in
Section~\ref{sec:lens}. Their move is to \emph{internalise the base}: $\Set$ is replaced by
$\mathbf{Loc}$ (or $\mathbf{Top}$), and small categories with retrofunctors become their localic
counterpart. Our $\DCont\cong\Cof$ is the $\Set$-level, detopologised ground case of the object
over which their adjunction sits; conversely their construction is a data point --- not yet an
answer --- for the question of what the morphism dictionary becomes over a base other than $\Set$.

Fairbanks, Carlson, and Spivak \cite{FCS26} characterise topological spaces as density comonads
of diagrams of subsets, explicitly motivated by the Ahman--Uustalu characterisation of categories
as the polynomial comonads on $\Set$; the resulting double category of continuous maps
``recovers the double category of functors and retrofunctors of Clarke and Di~Meglio'' --- the
same functor/cofunctor pairing that Corollary~\ref{cor:lens} reads as Get and Put. They cite
Ahman--Uustalu as settled prior art and do not touch the directed-container axioms. It is the
closest neighbour in vocabulary, and supplies context rather than a competing statement of the
morphism dictionary.

\section{Conclusion and outlook}\label{sec:conclusion}

The object dictionary --- directed containers are small categories --- extends to
morphisms, but not in the obvious way. A morphism of directed containers is a
\emph{cofunctor}, not a functor, because its position component runs contravariantly;
the precise statement is an isomorphism of categories $\DCont\cong\Cof$. Functors are
the mirror-image, covariant, notion, and are not morphisms of containers; the dictionary
does not extend to $\Cat$, failing already on a one-object source. The contravariant
variance is exactly the Put of a delta lens, so a directed-container morphism is the
update-propagating half of a bidirectional transformation.

The practical consequence is a correction to make wherever the chain is used: a ``map of
directed containers'' carries update (Put) semantics and must be composed as a cofunctor.
In supply-chain composition and agent orchestration --- where a map between two staged
structures should propagate a downstream change back upstream --- this is the difference
between the right morphisms and a silently different set of them, of even a different
cardinality.

Two directions remain. First, formalisation: $\Cof$ and Theorem~\ref{thm:main} are
natural targets for a Lean development sitting on the directed-container/category
dictionary, which would let the correction be enforced rather than remembered. Second,
and more speculatively, the lens reading suggests a measure of how much a map
``rewrites'' as it propagates an update; whether such a diversity invariant of a
directed-container morphism is itself compositional is a question we leave open, as it
lies beyond the bookkeeping settled here.

\begin{thebibliography}{99}

\bibitem{Aguiar}
M.~Aguiar.
\emph{Internal categories and quantum groups}.
PhD thesis, Cornell University, 1997.

\bibitem{ACU14}
D.~Ahman, J.~Chapman, and T.~Uustalu.
\emph{When is a container a comonad?}
Logical Methods in Computer Science \textbf{10}(3), 2014.

\bibitem{AU16}
D.~Ahman and T.~Uustalu.
\emph{Directed containers as categories}.
In \emph{Proc.\ MSFP 2016}, EPTCS \textbf{207}, 2016, pp.\ 89--98.

\bibitem{Clarke20}
B.~Clarke.
\emph{Internal lenses as functors and cofunctors}.
In \emph{Proc.\ ACT 2019}, EPTCS \textbf{323}, 2020, pp.\ 183--195.
\href{https://arxiv.org/abs/2009.06835}{arXiv:2009.06835}.

\bibitem{Clarke22}
B.~Clarke.
\emph{Delta lenses as coalgebras for a comonad}.
\href{https://arxiv.org/abs/2108.00390}{arXiv:2108.00390}, 2022.

\bibitem{JR13}
M.~Johnson and R.~Rosebrugh.
\emph{Delta lenses and opfibrations}.
Electronic Communications of the EASST \textbf{57}, 2013.

\bibitem{GRW26}
R.~Garner, A.~Renata, and N.~Wu.
\emph{Stone duality for monads}.
In \emph{Proc.\ MFPS 2026}, preliminary proceedings.
\href{https://arxiv.org/abs/2603.25710}{arXiv:2603.25710}.

\bibitem{FCS26}
A.~Fairbanks, Carlson, and D.~I.~Spivak.
\emph{Comonads as spaces}.
\href{https://arxiv.org/abs/2607.15091}{arXiv:2607.15091}, 2026.

\bibitem{SS}
B.~Shapiro and D.~I.~Spivak.
\href{https://arxiv.org/abs/2405.13157}{arXiv:2405.13157}, 2024;
on $\mathbf{Cat}^\sharp$, the category of small categories and cofunctors realised
inside $(\Poly,\triangleleft)$.

\end{thebibliography}

\end{document}
